\section{The Misconception of Egalitarianism}

\begin{quotex}
Every doctrine is true in what it asserts but false in, or because of, what it excludes.
\flright{\textsc{Leibniz}}
\end{quotex}


I've been looking at some writings from the ``right'', or ``tradition'', or ``conservatism'' of various persuasions and I see that there many misconceptions, misunderstandings, and half-truths (i.e., wrong in what they exclude). This is the first in a series of posts that will discuss some of them.

A prevalent myth is the assertion that Christianity is inherently ``egalitarian''; this opinion is so common that it is never substantiated. We'll assume for the time being that social egalitarian is what is meant, that is, there are no differences in social distinctions. We wonder, too, where in a social hierarchy the opponents of egalitarianism expect to end up. But getting back to the main point, let's just quote what Christians themselves have written on this point. Consider the following quotes carefully; the opposite of egalitarianism is hierarchy and authority. This may not be to everyone's liking.

\begin{quotex}
As it is suitable that the institution of a city or a kingdom be made according to the model of the institution of the world, similarly it is necessary to draw from divine government the order (ratio) or Logos in Greek of the government of a city ... That is why authors of the Middle Ages could not imagine Christianity without an Emperor, just as they could not imagine the Universal Church without a pope. Because if the world is governed hierarchically, Christianity cannot be otherwise.


\flright{\textsc{Valentin Tomberg}}
\end{quotex}


\begin{quotex}
When once men recognise, both in private and in public life that Christ is King, society will at last receive the great blessings of real liberty, well ordered discipline, peace and harmony. Our Lord's regal office invests the human authority of princes and rulers with a religious significance; it ennobles the citizen's duty of obedience.


\flright{\textsc{Pope Pius XI (1925)}}
\end{quotex}


\begin{quotex}
No one doubts that all men are equal one to another, so far as regards their common origin and nature, or the last end which each one has to attain, or the rights and duties which are thence derived. But, as the abilities of all are not equal, as one differs from another in the powers of mind or body, and as there are very many dissimilarities of manner, disposition, and character, it is more repugnamt to reason to endeavor to confine all within the same measure, and to extend complete equality to the institutions of civil life.

There is an almost infinite dissimilarity of men, as parts of the whole. If they are to be all equal, and each is to follow his own will, the State will appear most deformed.


\flright{\textsc{Pope Leo XIII (1884)}}
\end{quotex}


\begin{quotex}
It is not true that all have equal rights in civil society. It is not true that there exists non lawful social hierarchy.


\flright{\textsc{Pope Pius XI (1937)}}
\end{quotex}


\begin{quotex}
This is how one arrives at the dictatorship of the proletariat, the class hostile to the hierarchical principle, which latter, however, is the reflection of the divine order. This is why the proletariat professes atheism.


\flright{\textsc{Valentin Tomberg}}
\end{quotex}


To recapitulate: social hierarchy is the reflexion of divine order. That is why opponents of the hierarchical order, from Thomas Paine to the sanscullottes to the Bosheviks profess atheism. Atheists, then, who oppose egalitarianism, can do so only on the basis of the will to power and not on any moral authority.

Now the popes concede that there is an equality of nature --- the essence of being human --- but insist that the inequality of accidental qualities is legitimate and just. If is what those web sites mean by egalitarianism, they would be correct. In that case, they insist that there is a difference in human nature itself; that is, different individuals, groups, or races differ essentially. It is not clear what they would mean by that, especially since they base their arguments on genetic grounds (that is, accidental qualities). If there are essential differences among humans, (that is, differences by nature) or beings that appear outwardly to be human, then they need to make this clear.


\flrightit{Posted on 2010-03-30 by Cologero}

\begin{center}* * *\end{center}

\begin{footnotesize}
\begin{sffamily}
\texttt{Ben on 2011-06-25 at 15:57 said:}

I really want to accept what you've asserted here in this post, as I'm a supporter of conservative, traditional monarchy and its relation to Catholicism... But I just can't swallow the idea of essential racial inequality. It just has such a ring of 19th century social Darwinism that I can't stomach. Is there any way you could clarify this idea?

\hfill

\texttt{Cologero on 2011-06-25 at 19:37 said:}

Against my better judgment, Ben, I decided to allow your comment, as you seem to be a troll. First of all, we would encourage you to use your intellect in making judgments, not your dysphagia which is of no interest to anyone. There is nothing in that post that defends ``essential racial inequality'', and absolutely nothing at all that defends Darwinism. Actually, just the opposite, or did you forget to read it to the end?

\hfill

\texttt{Ben on 2011-06-25 at 23:13 said:}

No, I have no interest in trolling at all. I'm just trying to understand this political philosophy, which I certainly want to support. And I didn't mean to accuse you or your sources of believing what I alluded to earlier; my problem was that if ``the inequality of accidental qualities is legitimate and just,'' and genetics are accidental qualities, would that not imply that racial inequalities are legitimate?

\hfill

\texttt{Logres on 2011-06-26 at 01:27 said:}

Ben, your argument would have to assume as another premise that ``genetic material'' compromises the sum total of a human being. Actually, from another point of view starting from your very premises, you are the one defending what you are attacking. Humans have to ``be'' in a particular way, so genetic material is essential without being completely exhaustive.

\hfill

\texttt{Ben on 2011-06-26 at 11:13 said:}

So, in other words, there are other accidental qualities aside from what is handed down genetically? And in addition to that, there are the things Pope Leo XIII talked about as being common to humanity as a whole... So that a human being is a result of numerous diverse factors which contribute to his place in the hierarchy, yet even with the existence of the hierarchy there is still much that unites men.

\hfill

\texttt{Cologero on 2011-06-26 at 23:00 said:}

Ben, your obsession with race can only be the result of indoctrination by the American educational system.

The Traditional view is that there is a cosmic order, which is just; piety, that is, the duty of following that order, is just. This is opposed to the view of the gnostics who regard the cosmos as the creation of an evil demi-urge; hence, they support revolution and the overturning of every manifestation of cosmic order.

Find the post on Providence, Will and Destiny. The world process (destiny) is neither just nor unjust; an apple falling to the ground acts neither justly nor unjustly. Similarly, a genetic process is just what it is. We accept our fate (amor fati), this it he second trial; this is not a weakness, despite what you may read elsewhere. At a deeper level, we understand this as a privation of our own nature.

But the relations between men are a different story; they can be just or unjust, in conformance with the order of things or a violation of it. Injustices arising that way are not legitimate. But this opens the issue of what are one's duties? A man is obligated to rectify his own injustices, but is he obligated to rectify natural inequalities or those brought on by another man himself? The modern man, who is a gnostic at heart, will answer yes to the latter.

Then there are alleged inequalities of Providence. These are related to the ideas of Predestination and Predilection, which can't be addressed here. If you are Catholic, check out Garrigou-Lagrange's book Predestination.

\hfill

\texttt{Ben on 2011-06-27 at 14:47 said:}

It's amusing to me that you are able to see to the heart of my difficulty -- yes, I am unfortunately the product (to a certain extent) of the American public school system. My earliest education was in desegregated south Louisiana. As to my whether I'm Catholic, I'm in the process of conversion.

\hfill

\texttt{Sullen Eyes on 2023-03-29 at 22:47 said:}

It's been quite a journey for me to grapple with egalitarianism. I came to Guenon, Schuon, and Evola through the contemporary Internet-based far right, and there's a lot of baggage that comes with that. It's been a long journey to unlearn that perspective at the reflexive, emotional level.

I've primarily based my practice around Vaishnava bhakti meditation, and the conception of ``egalitarianism'' in that tradition is very different from the contemporary political conception (if egalitarian is even a proper word to use for it). As I understand it: one's obligation to love God applies to all of His creation, even if there are differences of social and spiritual rank between beings as well as differences in orientation and vocation. Furthermore, the detachment from egoistic concerns naturally causes a sort of impartiality among sages.

I also find it somewhat interesting how little attention Sikhism, with it's great emphasis on bhakti-based egalitarianism, has received among western writers whether they be writing about tradition or some sort of political philosophy. Reading the Japji Sahib has been illuminating for one who is as imbalanced as me.

\end{sffamily}
\end{footnotesize}

